% Part: counterfactuals
% Chapter: introduction
% Section: paradoxes-material

\documentclass[../../../include/open-logic-section]{subfiles}

\begin{document}

\olfileid[hi]{cnt}{int}{par}

\olsection{भौतिक सशर्त के विरोधाभास}

अंग्रेज़ी के ``यदि \dots तो \dots'' कथनों को संकेतित करने के लिए
$!A \lif !B$ के प्रयोग की आलोचना करने वालों में सी.~आई. लुईस सबसे पहले
लोगों में थे। लुईस व्हाइटहेड और रसेल की \emph{प्रिंसिपिया मैथमेटिका} में
भौतिक सशर्त के प्रयोग की आलोचना कर रहे थे, जहाँ $\lif$ को ``निहित करता है''
पढ़ा गया था। लुईस ने उचित ही शिकायत की कि यदि $\lif$ का अर्थ ``निहित करता
है'' है, तो कोई भी असत्य प्रतिज्ञप्ति~$p$ यह निहित करती है कि $p$,~$q$ को
निहित करता है, क्योंकि~$p$ के असत्य होने पर $p \lif (p \lif q)$ सत्य है;
और कोई भी सत्य प्रतिज्ञप्ति~$q$ यह निहित करती है कि $p$,~$q$ को निहित करता
है, क्योंकि~$q$ के सत्य होने पर $q \lif (p \lif q)$ सत्य है।

तर्कशास्त्री निश्चय ही जानते हैं कि \emph{निहितार्थ}, अर्थात तार्किक
निष्कर्षण, कोई संयोजक नहीं बल्कि !!{formula}ों या कथनों के बीच संबंध है। अतः
भ्रम से बचने के लिए हमें $\lif$ को ``निहित करता है'' नहीं पढ़ना चाहिए।
\footnote{गणितज्ञ और कंप्यूटर वैज्ञानिक अब भी ``$\lif$'' को ``निहित करता
  है'' व्यापक रूप से पढ़ते हैं, यद्यपि दार्शनिक उसे ``केवल तभी'' पढ़कर लुईस
  द्वारा रेखांकित भ्रमों से बचने का प्रयास करते हैं।} जब तक हम ऐसा नहीं
पढ़ते, लुईस की विशेष चिंता उत्पन्न ही नहीं होती: भले ही $p$ को किसी ऐसे
अंग्रेज़ी वाक्य का संकेत मानें जो संयोगतः असत्य है, फिर भी $p$,~$q$ को
``निहित'' नहीं करता। $p \Entails q$ है या नहीं यह निर्धारित करने के लिए
हमें \emph{सभी} !!{valuation}ों पर विचार करना होगा; और किसी संयोगतः असत्य
वाक्य को संकेतित करने के लिए~$p$ का उपयोग करने पर भी $p \Entails/ q$।

फिर भी लुईस के निम्न जैसे ``यदि \dots तो \dots'' कथनों में कुछ विचित्र है:
\begin{quote}
यदि चंद्रमा हरे पनीर का बना है, तो $2+2=4$।
\end{quote}
और निम्न अनुमानों में भी:
\begin{quote}
  चंद्रमा हरे पनीर का बना नहीं है। अतः यदि चंद्रमा हरे पनीर का बना है, तो
  $2+2=4$।

  $2+2 = 4$। अतः यदि चंद्रमा हरे पनीर का बना है, तो $2+2=4$।
\end{quote}
फिर भी यदि ``यदि \dots तो \dots'' केवल $\lif$ होता, तो वाक्य निर्विवाद रूप
से सत्य और अनुमान निर्विवाद रूप से मान्य होते।

एक अन्य उदाहरण सर्वसत्यता $(!A \lif !B) \lor (!B \lif !A)$ से संबंधित है।
इससे संकेत मिलता है कि यदि समाचारपत्र से दो विधानवाचक वाक्य~$S$ और $T$
यादृच्छिक रूप से लें, तो ``यदि $S$ तो $T$, या यदि $T$ तो~$S$'' वाक्य सत्य
होना चाहिए।

\end{document}
